% ============================================================
% SEÇÃO 3 – ESPECIFICAÇÃO DOS CASOS DE USO
% ============================================================

\section{ESPECIFICAÇÃO DOS CASOS DE USO}

\subsection{Login e visualização do mapa}

\begin{longtable}{|p{4cm}|p{10cm}|}
  \hline
  \textbf{Nome} & Login e visualização do mapa \\
  \hline
  \textbf{Descrição} & Usuário com conta cadastrada consegue iniciar sessão no sistema e visualizar no mapa os eventos, ocorrências ou informações atualizadas sobre o que está acontecendo na cidade. \\
  \hline
  \textbf{Ator} & Usuário do aplicativo \\
  \hline
  \textbf{Pré-condições} &
    \begin{itemize}
      \item Usuário precisa possuir conta cadastrada.
      \item Sistema deve estar conectado à internet para carregar informações do mapa.
    \end{itemize} \\
  \hline
  \textbf{Pós-condições} &
    \begin{itemize}
      \item Usuário consegue acessar o aplicativo.
      \item Usuário visualiza o mapa com os eventos/ocorrências ativas na cidade.
    \end{itemize} \\
  \hline
  \textbf{Requisitos funcionais} & RF002 — Login do usuário \newline RF007 — Exibição de mapa interativo \newline RF008 — Filtrar por categoria \\
  \hline
  \textbf{Fluxo básico} &
    \begin{enumerate}
      \item Sistema abre a tela de iniciar sessão com campos a serem preenchidos.
      \item Usuário preenche corretamente os campos de login.
      \item Usuário aperta o botão para iniciar sessão.
      \item Sistema verifica se os dados são válidos.
      \item Sistema inicia a sessão, retornando tokens JWT (access + refresh).
      \item Sistema redireciona o usuário para a tela principal com o mapa da cidade.
      \item Sistema carrega e exibe no mapa os eventos e ocorrências.
      \item Usuário visualiza no mapa o que está acontecendo na cidade.
    \end{enumerate} \\
  \hline
  \textbf{Fluxo alternativo} & \textbf{FA01 — Dados inválidos no login} \newline
    Sistema exibe mensagem de erro referente aos dados inválidos e permite nova tentativa. \\
  \hline
\end{longtable}

\subsection{Usuário busca e encontra o evento que deseja participar}

\begin{longtable}{|p{4cm}|p{10cm}|}
  \hline
  \textbf{Nome} & Busca de evento próximo no mapa \\
  \hline
  \textbf{Descrição} & Usuário autenticado acessa o mapa, filtra eventos próximos e visualiza as informações detalhadas de um evento para avaliar sua participação. \\
  \hline
  \textbf{Ator} & Usuário do aplicativo \\
  \hline
  \textbf{Pré-condições} &
    \begin{itemize}
      \item Usuário precisa estar logado no aplicativo.
      \item Sistema deve ter acesso à base de eventos cadastrados.
      \item Sistema deve estar conectado à internet.
    \end{itemize} \\
  \hline
  \textbf{Pós-condições} &
    \begin{itemize}
      \item Usuário visualiza no mapa os eventos próximos.
      \item Usuário consegue visualizar os detalhes de um evento para decidir se participa.
    \end{itemize} \\
  \hline
  \textbf{Requisitos funcionais} & RF007 — Visualizar mapa \newline RF008 — Filtrar por categoria \newline RF005 — Registrar evento \\
  \hline
  \textbf{Fluxo básico} &
    \begin{enumerate}
      \item Usuário seleciona a aba ``Eventos'' no mapa.
      \item Sistema exibe no mapa marcadores representando os eventos cadastrados.
      \item Usuário pode filtrar por categoria (Festas, Esportes, Educação, etc.).
      \item Usuário toca em um marcador de evento no mapa.
      \item Sistema abre card com as informações do evento (nome, endereço, horário, categoria, descrição).
      \item Usuário visualiza os detalhes e avalia se deseja participar.
    \end{enumerate} \\
  \hline
  \textbf{Fluxo alternativo} & \textbf{FA01 — Nenhum evento encontrado} \newline
    Sistema exibe o mapa sem marcadores. \newline\newline
    \textbf{FA02 — Falha de conexão} \newline
    Sistema exibe mensagem de erro de rede. \\
  \hline
\end{longtable}

\subsection{Usuário denuncia problema em sua rua}

\begin{longtable}{|p{4cm}|p{10cm}|}
  \hline
  \textbf{Nome} & Denúncia feita no sistema \\
  \hline
  \textbf{Descrição} & Usuário acessa o aplicativo, registra uma denúncia informando a existência de um problema em sua rua e envia a localização e dados relevantes para que o órgão responsável possa avaliar e tomar providências. \\
  \hline
  \textbf{Ator} & Usuário \\
  \hline
  \textbf{Pré-condições} &
    \begin{itemize}
      \item Usuário precisa estar logado no aplicativo.
      \item Usuário precisa estar em área com acesso à internet.
      \item Sistema deve ter acesso ao serviço de localização (GPS) ou permitir seleção manual.
    \end{itemize} \\
  \hline
  \textbf{Pós-condições} &
    \begin{itemize}
      \item Denúncia é registrada no sistema e exibida no mapa.
      \item Usuário recebe confirmação de envio.
    \end{itemize} \\
  \hline
  \textbf{Requisitos funcionais} & RF003 — Registrar reclamação \newline RF007 — Visualizar mapa \\
  \hline
  \textbf{Fluxo básico} &
    \begin{enumerate}
      \item Usuário acessa o aplicativo e toca no botão ``+'' para criar publicação.
      \item Sistema abre a tela de registro de reclamação.
      \item Usuário seleciona o tipo de problema (Infraestrutura, Segurança, Limpeza, Trânsito, Outros).
      \item Usuário informa o endereço ou usa GPS automático.
      \item Usuário preenche a descrição do problema.
      \item Usuário adiciona foto (opcional).
      \item Usuário informa a data de ocorrência.
      \item Usuário pressiona ``Criar Reclamação''.
      \item Sistema valida e registra a denúncia no banco de dados.
      \item Sistema retorna feedback positivo ao usuário.
    \end{enumerate} \\
  \hline
  \textbf{Fluxo alternativo} & \textbf{FA01 — Falha ao obter localização via GPS} \newline
    Sistema oferece opção para o usuário informar o endereço manualmente. \newline\newline
    \textbf{FA02 — Campos obrigatórios incompletos} \newline
    Sistema exibe mensagem de erro e mantém o formulário para preenchimento. \\
  \hline
\end{longtable}
